\section{Experiments}
\label{sec:experiments}

\subsection{Setup}
\label{sec:exp-setup}

For development we used one calibrated multi-camera capture of a masked studio subject
with 26 cameras at $5328 \times 4608$, fitted compact captures for all views and split
them into 19 fitting views and 7 compact validation views.  This scene is outcome-exposed,
since it drove the design decisions, and is thus reported as development evidence only.
\todo{Add the confirmatory data.  At least one fresh calibrated dome scene with untouched
test cameras, plus Mip-NeRF~360~\citep{barron2022_mipnerf360} and Tanks and
Temples~\citep{knapitsch2017_tanks} scenes with the capture stage run at benchmark
resolution.  Freeze the train, validation and test protocol per scene before any
confirmatory run and touch the test cameras exactly once for the final tables, see T-4
and T-5.}

We report held-out PSNR, SSIM~\citep{wang2004_ssim} and LPIPS~\citep{zhang2018_lpips} in
every cell, alpha IoU and exterior leakage for masked subjects, the resource metrics of
\cref{sec:protocol} and the compact-domain risks $\Ju$ and $\Jq$ for diagnostics.
Geometry error against a reference depth or mesh is reported wherever a reference exists,
since PSNR is blind to floaters and floaters are the failure mode initialization is
supposed to prevent.

\subsection{Arms and baselines}
\label{sec:exp-baselines}

\begin{table}[t]
  \centering\footnotesize
  \setlength{\tabcolsep}{4pt}
  \begin{tabular}{@{}lp{3.1cm}p{3.6cm}p{3.6cm}@{}}
    \toprule
    Arm & Supervision & Initialization & Status \\
    \midrule
    Original 3DGS            & RGB & COLMAP sparse
      & \tbd, needs COLMAP (T-7) \\
    Orig.\ 3DGS, compr.\ RGB & matched-byte JPEG or AVIF & COLMAP sparse
      & \tbd \\
    RGB control              & RGB & frozen direct-path init, byte-identical
      & \tbd \\
    Direct, random init      & compact fields & bounded random
      & \tbd, schedule to freeze (T-3) \\
    Direct, SfM init         & compact fields & COLMAP sparse
      & \tbd, needs COLMAP (T-7) \\
    Tomographic              & compact fields & Beam Fusion + refinement
      & \tbd \\
    \bottomrule
  \end{tabular}
  \caption{Arms of the comparison.  The RGB control isolates the supervision variable
    through an identical initialization, and the Original-3DGS rows anchor against the
    community-standard pipeline.  \todo{All quantitative cells are pending.  A
    fixed-topology RGB control is added to separate supervision value from capacity
    growth.}}
  \label{tab:arms}
\end{table}

\subsection{Main results}
\label{sec:exp-main}

\begin{table}[t]
  \centering\footnotesize
  \setlength{\tabcolsep}{4pt}
  \begin{tabular}{@{}lcccccccc@{}}
    \toprule
    & \multicolumn{3}{c}{Held-out quality} & \multicolumn{3}{c}{Resources (peak)} & & \\
    \cmidrule(lr){2-4}\cmidrule(lr){5-7}
    Arm & PSNR$\uparrow$ & SSIM$\uparrow$ & LPIPS$\downarrow$
        & GPU mem & host RAM & disk & IO & time \\
    \midrule
    Original 3DGS            & \tbd & \tbd & \tbd & \tbd & \tbd & \tbd & \tbd & \tbd \\
    Orig.\ 3DGS, compr.\ RGB & \tbd & \tbd & \tbd & \tbd & \tbd & \tbd & \tbd & \tbd \\
    RGB control (same init)  & \tbd & \tbd & \tbd & \tbd & \tbd & \tbd & \tbd & \tbd \\
    Direct, random init      & \tbd & \tbd & \tbd & \tbd & \tbd & \tbd & \tbd & \tbd \\
    Direct, SfM init         & \tbd & \tbd & \tbd & \tbd & \tbd & \tbd & \tbd & \tbd \\
    Tomographic              & \tbd & \tbd & \tbd & \tbd & \tbd & \tbd & \tbd & \tbd \\
    \bottomrule
  \end{tabular}
  \caption{\redcaption{Main table with held-out quality and measured resource peaks per
    arm, as the median over at least three fresh-process repeats with the spread in the
    appendix and per-scene sub-tables.  Entirely pending the controlled experiment, see
    T-6.}}
  \label{tab:main-results}
\end{table}

\begin{figure}[t]
  \centering
  \figslot{6cm}{%
    \textbf{Figure 9, qualitative comparison.}  Two to three held-out test cameras per
    scene.  \textbf{Left:} RGB-control render.  \textbf{Middle:} direct-path render.
    \textbf{Right:} difference maps and crops of the hardest regions like hair, fabric
    edges and specular highlights.  The purpose is to make the parity or its failure
    inspectable.  Uses the final models of \cref{tab:main-results}.}
  \caption{\redcaption{Qualitative held-out comparison of RGB-backed and
    compact-supervised training.}}
  \label{fig:parity-grid}
\end{figure}

Two development results bound the expectations.  A full-scale development run on the
exposed scene, with all 26 views fitted and RGB-based optimization, reached $28.76$~dB
native foreground PSNR, $0.974$ SSIM and $0.0173$ LPIPS at 100k Gaussians after 128k
updates, which confirms that the capture and optimization machinery scales to native
resolution.\evidence{ara/logic/claims.md C29}  The compact-only carrier pipeline reached
compact validation risks $\Jq = 0.00422$ and $\Ju = 0.00786$ on the same scene under the
strict no-image boundary, which confirms that the compact objective optimizes stably end
to end.\evidence{ara/logic/claims.md C30, C31}  Neither run includes held-out cameras or a
resource comparison, so neither substitutes for \cref{tab:main-results}.

\subsection{Initialization comparison}
\label{sec:init-comparison}

The value of the tomographic initializer cannot be inferred from its own result alone.  We
test it by a matched comparison at exact count and compute with all downstream settings
identical.  The arms are the tomographic initialization with refinement, the image-free
placement control which scores rays by querying all fields, and calibration-bounded
random as the lower anchor.  All three feed the identical compact optimizer and evaluation
samples.  We report initialization quality, fixed intermediate checkpoints, the risk-curve
area, the selected endpoint and paired seed wins on all scenes.  The construction time and
peak memory of the initializer are accounted separately, since the initializer must pay
for itself end to end.  A converged tie after an initialization win means the initializer
improves startup but not final capacity, and the text will say exactly that if it
happens.

\begin{table}[t]
  \centering\small
  \begin{tabular}{@{}lccccc@{}}
    \toprule
    Initialization & init $\Jq$ & AUC($\Jq$) & final $\Jq$ & held-out PSNR & seed wins \\
    \midrule
    Tomographic + refinement & \tbd & \tbd & \tbd & \tbd & \tbd \\
    Field-query placement    & \tbd & \tbd & \tbd & \tbd & \tbd \\
    Bounded random           & \tbd & \tbd & \tbd & \tbd & \tbd \\
    \bottomrule
  \end{tabular}
  \caption{\redcaption{Matched initialization comparison per scene with paired seeds,
    including a companion column for initializer build time and peak memory.  Entirely
    pending, see T-8.}}
  \label{tab:beam-value}
\end{table}

\begin{figure}[t]
  \centering
  \figslot{5.5cm}{%
    \textbf{Figure 10, risk trajectories.}  Validation $\Jq$ on a logarithmic axis over
    optimizer updates for the three initializations of \cref{tab:beam-value}, seeds as
    thin lines, medians bold, with initialization points and the fixed checkpoints
    annotated.  A second panel shows held-out PSNR over wall clock, where the tomographic
    arm starts with an offset equal to its construction time, which makes the
    pay-for-itself question visible.  Uses the matched comparison runs.}
  \caption{\redcaption{Risk trajectories for the tomographic, placement and random
    initializations.}}
  \label{fig:time-to-quality}
\end{figure}

\todo{Freeze the primary endpoint, the materiality margin and the seed-win rule of this
comparison before the confirmatory runs, see T-9.  If the tomographic arm ties or loses,
the direct path remains the recommended pipeline, the outcome is reported in full and the
result of \cref{sec:memory} is unaffected, see \cref{sec:rules}.}

\subsection{Ablation study}
\label{sec:ablations}

\begin{table}[t]
  \centering\footnotesize
  \setlength{\tabcolsep}{4pt}
  \begin{tabular}{@{}lp{7.1cm}l@{}}
    \toprule
    Ablation & Development verdict (single scene) & Confirmatory \\
    \midrule
    Without covariance repair
      & repair lowers validation $\Jq$ by $63.2\%$, 3/3 paired wins \evidence{C30}
      & \tbd \\
    One-sided repair residual
      & shrink-biased, worsens $\Jq$ by $51.6\%$, replaced by \cref{eq:cov-residual}
        \evidence{C30}
      & \tbd \\
    One compact phase versus two
      & second phase reaches $0.717\times$ its stopping $\Jq$, 3/3 \evidence{C30}
      & \tbd \\
    Phase-one means frozen
      & fails the $2\%$ non-inferiority gate \evidence{C30}
      & \tbd \\
    Means-only optimization
      & $4.18\times$ the all-family $\Jq$ \evidence{C30}
      & \tbd \\
    Without containment prune
      & prune removes $5.3$ to $5.4\%$ at $2.7\%$ $\Jq$ cost, zero violations
        \evidence{C30, C31}
      & \tbd \\
    Prune before versus after phase two
      & before selected \evidence{C30}
      & \tbd \\
    Phase-two means trainable versus frozen
      & frozen passes at ratio $1.015$ $\Jq$ \evidence{C30}
      & \tbd \\
    Clone stage (density-preserving)
      & fails the $5\%$ materiality gate, no growth \evidence{C30}
      & \tbd \\
    SH0 versus incremental SH3
      & SH3 gains $2.3$ to $3.0\%$, below the $5\%$ floor \evidence{C30}
      & \tbd \\
    Soft support barrier \cref{eq:soft-support}
      & \todo{freeze, then test}
      & \tbd \\
    \bottomrule
  \end{tabular}
  \caption{Ablation grid for the tomographic path.  Development verdicts are frozen
    paired-root results on one outcome-exposed scene with 19 fitting views and compact
    validation.  Evidence keys refer to rows of \repopath{ara/logic/claims.md}.
    \todo{The confirmatory column with fresh scenes, held-out cameras and at least three
    paired seeds under preregistered gates is entirely pending, see T-10.}}
  \label{tab:ablations}
\end{table}

The development column of \cref{tab:ablations} closes the implementation-policy question
on one scene, i.e.\ which stages the refinement runs and in which order.  It does not
close the value question of \cref{tab:beam-value}, which is a comparison against the
non-tomographic arms, and it does not establish cross-scene generalization.

\subsection{Diagnostics}
\label{sec:diagnostics}

Per initializer run we report the contributor-count histogram, per-view uniqueness, gate
survival rates, covariance eigenvalue spectra and lineage hashes.  Per refinement we
report carrier survival, per-family optimizer motion, the covariance residual
distribution of \cref{eq:cov-residual}, containment violation counts before and after the
prune and the fixed-topology recheck.  These quantities make a false-association failure,
the designed risk of \cref{sec:ci-fusion}, visible instead of silent.

\begin{figure}[t]
  \centering
  \figslot{5cm}{%
    \textbf{Figure 11, diagnostics panel.}  Four panels.  \textbf{Top left:} the
    contributor-count histogram.  \textbf{Top right:} carrier survival across stages as
    bars.  \textbf{Bottom left:} the distribution of covariance residuals
    \cref{eq:cov-residual} before and after repair.  \textbf{Bottom right:} containment
    violations per view before the prune.  Uses the diagnostics dictionaries already
    emitted by the implementation on the development scene under \repopath{runs/}.}
  \caption{\redcaption{Initializer and refinement diagnostics.  Producible from existing
    run artifacts.}}
  \label{fig:beam-diagnostics}
\end{figure}
